% Part: first-order-logic
% Chapter: tableaux
% Section: identity

\documentclass[../../../include/open-logic-section]{subfiles}

\begin{document}

\olfileid{fol}{tab}{ide}

\olsection{\usetoken{P}{tableau} با \usetoken{S}{identity}}

!!^{tableau}s دارای !!{identity} به قواعد استنتاج بیشتری نیاز دارند.
قواعد $\eq$ چنین‌اند ($t$، $t_1$ و $t_2$ ترم‌های بسته‌اند):

\begin{defish}
\AxiomC{}
\RightLabel{$\eq$}
\UnaryInfC{\sFmla{\True}{\eq[t][t]}}
\DisplayProof
\hfill
\AxiomC{\sFmla{\True}{\eq[t_1][t_2]}}
\noLine
\UnaryInfC{\sFmla{\True}{!A(t_1)}}
\RightLabel{$\TRule{\True}{\eq}$}
\UnaryInfC{\sFmla{\True}{!A(t_2)}}
\DisplayProof
\hfill
\AxiomC{\sFmla{\True}{\eq[t_1][t_2]}}
\noLine
\UnaryInfC{\sFmla{\False}{!A(t_1)}}
\RightLabel{$\TRule{\False}{\eq}$}
\UnaryInfC{\sFmla{\False}{!A(t_2)}}
\DisplayProof
\end{defish}
توجه کنید که برخلاف همهٔ قواعد دیگر، $\TRule{\True}{\eq}$ و
$\TRule{\False}{\eq}$ مستلزم آن‌اند که \emph{دو} !!{formula}s علامت‌دار
از پیش روی شاخه ظاهر شده باشند، یعنی هم
$\sFmla{\True}{\eq[t_1][t_2]}$ و هم $\sFmla{S}{!A(t_1)}$.

\begin{ex}
اگر $s$ و $t$ ترم‌های بسته باشند، آنگاه
$\eq[s][t], !A(s) \Proves !A(t)$:
\begin{oltableau}
  [\sFmla{\False}{\formula{A}(t)}, just = \TAss
    [\sFmla{\True}{\eq[s][t]}, just = \TAss
      [\sFmla{\True}{\formula{A}(s)}, just = \TAss
        [\sFmla{\True}{\formula{A}(t)}, just={\TRule{\True}{\eq}[2, 3]}, close]
      ]
    ]
  ]
\end{oltableau}
این اصل را شاید با نام اصل جانشینی این‌همان‌ها، یا قانون لایبنیتس،
بشناسید.

!!^{tableau}s اثبات می‌کنند که $\eq$ متقارن است؛ یعنی
$\eq[s_1][s_2] \Proves \eq[s_2][s_1]$:
\begin{oltableau}
  [\sFmla{\False}{\eq[s_2][s_1]}, just = \TAss
    [\sFmla{\True}{\eq[s_1][s_2]}, just = \TAss
      [\sFmla{\True}{\eq[s_1][s_1]}, just = {$\eq$}
        [\sFmla{\True}{\eq[s_2][s_1]}, just = {\TRule{\True}{\eq}[2, 3]}, close]
      ]
    ]
  ]
\end{oltableau}
در اینجا، سطر $2$ نخستین پیش‌شرط
!!{formula}~$\sFmla{\True}{\eq[s_1][s_2]}$ قاعدهٔ
$\TRule{\True}{\eq}$ است. سطر~$3$ پیش‌شرط دوم
به‌شکل $\sFmla{\True}{!A(s_1)}$ است---$!A(x)$ را
$\eq[x][s_1]$ در نظر بگیرید؛ آنگاه $!A(s_1)$ همان
$\eq[s_1][s_1]$ و $!A(s_2)$ همان $\eq[s_2][s_1]$ است.

این تابلوها همچنین اثبات می‌کنند که $\eq$ متعدی است؛ یعنی
$\eq[s_1][s_2], \eq[s_2][s_3] \Proves \eq[s_1][s_3]$:
\begin{oltableau}
  [\sFmla{\False}{\eq[s_1][s_3]}, just = \TAss
    [\sFmla{\True}{\eq[s_1][s_2]}, just = \TAss
      [\sFmla{\True}{\eq[s_2][s_3]}, just = \TAss
        [\sFmla{\True}{\eq[s_1][s_3]}, just = {\TRule{\True}{\eq}[3, 2]}, close]
      ]
    ]
  ]
\end{oltableau}
در این !!{tableau}، نخستین پیش‌شرط !!{formula} قاعدهٔ
$\TRule{\True}{\eq}$ سطر~$3$، یعنی
$\sFmla{\True}{\eq[s_2][s_3]}$ است ($s_2$ نقش~$t_1$ و $s_3$
نقش~$t_2$ را دارد). پیش‌شرط دوم، به‌شکل
$\sFmla{\True}{!A(s_2)}$، سطر~$2$ است. در اینجا $!A(x)$ را
$\eq[s_1][x]$ در نظر بگیرید؛ در این صورت $!A(s_2)$ به
$\eq[s_1][s_2]$ (یعنی سطر~$2$) و $!A(s_3)$ به
!!{formula}~$\eq[s_1][s_3]$ در نتیجه بدل می‌شود.
\end{ex}

\begin{prob}
برای موارد زیر !!{tableau}s بسته ارائه کنید:
\begin{enumerate}
\item $\sFmla{\False}{\lforall[x][\lforall[y][((x = y \land !A(x))
      \lif !A(y))]]}$
\item $\sFmla{\False}{\lexists[x][(!A(x) \land
    \lforall[y][(!A(y) \lif y = x)])]}$,\\
  $\sFmla{\True}{\lexists[x][!A(x)] \land
    \lforall[y][\lforall[z][((!A(y) \land !A(z)) \lif y = z)]]}$
\end{enumerate}
\end{prob}

\end{document}
